% ==================== 章节2：大语言模型的前世今生 ====================

% --- 章节标题页（竖排） ---
\hirochapter{大语言模型的前世今生}

% --- 小剧场对话 ---
\section*{小剧场}
\phantomsection
\addcontentsline{toc}{section}{小剧场}

午后的初星学园很安静。练习室的窗帘被风吹起一点。广站在白板前，看见上一节课没有擦干净的粉痕。

Producer 把马克笔递给她，没有马上松手。

“站得住吗？”

“三分钟。”我看了一眼墙上的钟，“如果中途坐下，就把它当作实验误差。”

“那就两分钟半。”

“嗯。很过分，だ、ね。”

广这样说着，还是接过了笔。

“今天讲大语言模型？”

她点头。接过笔，先没有写标题，而是在白板中央写下一行：

\begin{quote}
给定前面的文字，预测后面会出现什么。
\end{quote}

“只是这样？”

“只是这样。”广说，“所以很有趣。”

她把“预测”两个字圈起来。

“如果预测得很差，就是胡言乱语。如果预测得好一点，就是自动补全。如果预测得非常好，就会像是在回答问题、写文章、解释代码。”

Producer 看着白板：“那它是真的理解了吗？”

广停了一下，指尖扶住白板边缘，确认自己还算站得稳定。

“这个问题太快了。请再准确一点。”

她在白板左边写下“规则”，中间写下“统计”，右边写下“规模”。

“先看它怎么出生。一个东西从很笨，到很有用，中间不一定有某个瞬间突然获得灵魂。很多时候，它只是把同一件事做到了极限附近。”

“预测下一个词？”

“嗯。”

她把马克笔放回笔槽。

“做不到。再预测一次。”

她轻轻笑了一下。

“这里就是今天的课题。”

\clearpage

% --- 课堂讲义正文（IEEE风格，单栏） ---
\hiropapertitle{大语言模型的前世今生}{広}{初星学园}

\begin{abstract}
大语言模型的核心训练目标可以被概括为：给定上下文，预测下一个 token。这个目标很小，却在一次次撞到语言的边界之后，被迫长出新的结构：规则写不完真实语言，统计模型看不远，循环记忆训练太慢，预训练模型会续写却不一定会帮忙，聊天助手会帮忙却不一定可靠，能调用工具以后又必须面对权限、审计与责任。这堂“筱澤広小课堂”沿着这条问题链，结合代表性论文、公开产品节点、少量公式与原创示意图，说明大语言模型如何从“会接话”发展到“能协作”，并讨论它的能力边界、幻觉风险、系统可靠性与 agent 化趋势。
\end{abstract}

\keywords{大语言模型；Transformer；预训练；指令微调；RLHF；RAG；推理模型；AI agent}

\section{太小的起点：预测下一个 token}

如果只看结果，大语言模型很像突然出现的东西。

昨天我们还在搜索框里输入关键词，今天就可以对一个聊天窗口说：“帮我解释这段代码”“把会议纪要整理成待办”“用更温柔的语气改写这封邮件”。它会回答，会追问，会犯错，也会在某些时刻显得异常聪明。

于是问题来了：它到底是什么？

最短的答案是：大语言模型是一种被大规模训练过的概率模型。它阅读上下文，然后预测下一个 token。所谓 token，可以先粗略理解成文字被切成的小片段。一个汉字、一个英文词、一段词根、一个标点，都可能成为 token。

这听起来像文字接龙。可是，文字接龙接到极高水平时，会碰到很多我们原本以为属于“理解”的事情。要接好一句法律解释，它要知道术语关系；要接好一段代码，它要知道语法和运行结果；要接好一封道歉邮件，它要知道语气、责任和人际边界。

所以，广没有急着判断它有没有“真正理解”。她先把问题收窄：

\begin{quote}
为什么预测下一个 token 这件事，可以一路长成今天的大语言模型？
\end{quote}

这不是一条平滑的直线。更像一串很小的失败：规则覆盖不了，于是需要统计；统计看不远，于是需要表示和记忆；记忆会累、训练会慢，于是需要注意力；模型会续写，却不一定听题；模型会听题，却不一定说真话；模型能调用工具以后，又必须被边界约束。

也就是说，这堂课不是把技术名词排成年表。白板上只留下一个问题：

\begin{quote}
每一代方法解决了什么，又把下一个问题暴露在哪里？
\end{quote}

\begin{figure}[htbp]
\centering
\begin{tikzpicture}[>=stealth, every node/.style={font=\small}]
  \node[draw, rounded corners, align=center, text width=2.8cm] (ctx) at (0,0) {上下文\\今天的课题是};
  \node[draw, rounded corners, align=center, text width=2.8cm] (model) at (4,0) {语言模型\\计算概率分布};
  \node[draw, rounded corners, align=center, text width=2.8cm] (next) at (8,0) {下一个 token\\``大''、``AI''、``...''};
  \draw[->] (ctx) -- (model);
  \draw[->] (model) -- (next);
\end{tikzpicture}
\caption{语言模型的最小直觉：读上下文，预测下一个 token。}
\end{figure}

这张图就是课堂的起点。后面的复杂技术，都会回到它身上：怎样表示上下文，怎样计算概率，怎样训练参数，怎样让模型更听人话，怎样让它查资料、用工具、减少胡编。

不过，“预测”这个词本身也容易让人误会。它不是简单地从词典里挑一个最常见的词，也不是像搜索引擎那样把某个网页片段搬出来。更接近的说法是：模型把前文压缩成一组数值表示，再在整个词表上计算一个概率分布。这个分布里，某些 token 的概率高，某些 token 的概率低。生成时，模型根据这个分布选择一个 token，把它接到上下文后面，再重复下一轮。

Producer 举手：“token 是字吗？还是词？”

“都可能。”广说，“也可能是半个词、一个标点，或者一段常见字符。”

tokenizer 会先把原始文本切成模型词表里的小单位。中文里，一个 token 有时像一个字，有时像一个词；英文里，一个 token 可能是完整单词，也可能是词根或后缀。模型真正处理的不是纸面上的“句子”，而是一串 token 编号。编号查表变成向量，向量经过很多层计算，最后再变成下一个 token 的概率。

她在“切分”下面画了一条线。

“第一步不是理解。”

“第一步是把文字变成机器能计算的颗粒。”

所以，一次回答其实是一串很多小决定。每个 token 都只是一步，几百上千步连起来，才变成一段解释、一封邮件或一段代码。也正因为它是一步一步生成的，前面一步的小偏差会影响后面的路线。模型看起来是在“一口气回答”，内部却更像在一格一格铺路。

这也是为什么同一个问题有时会得到不同答案。模型并不是只保存了一个唯一答案，而是在一个概率空间里行动。温度、采样策略、上下文细节，都会改变它最后走到哪里。先理解这一点，后面的“幻觉”“推理”“工具调用”才不会显得神秘。

广把这里讲得很慢。

“先不要把它想成会说话的人。”

她停了一下，把马克笔换到左手，补上第二句。

“先把它想成一台非常认真、非常庞大的预测机器。”

\subsection*{数学骨架：把“接话”写成概率}

论文里讲语言模型，通常不会从“会聊天”开始，而是从一个概率分解开始。假设一段文本被切成 token 序列：

\[
x_1,x_2,\ldots,x_T
\]

语言模型要估计整段文本出现的概率。按照链式法则，可以写成：

\[
P(x_{1:T})=\prod_{t=1}^{T}P(x_t\mid x_{<t})
\]

这里的 \(x_{<t}\) 表示第 \(t\) 个 token 之前的所有 token。也就是说，模型并不是一次性吐出整篇文章，而是在每个位置上问自己：已经看到前面这些内容，下一个 token 应该是什么？

训练时，模型会在大量文本上反复猜答案。猜错了，就调整参数。常见的训练目标可以写成负对数似然：

\[
\mathcal{L}=-\sum_{t=1}^{T}\log P(x_t\mid x_{<t})
\]

这个式子不用吓人。它的意思很朴素：真实答案出现时，模型给它的概率越高，损失越小；模型越不相信真实答案，损失越大。训练就是让这个损失下降。

Producer 在旁边写下“老师批改”。

“差不多。”广说，“但这里没有老师逐句写评语。答案藏在文本自己里面。”

这就是自监督学习的直觉。对于一段现成文本，前面的 token 可以自动变成题目，真实出现的下一个 token 可以自动变成答案。模型不需要人类给每一句话贴标签，也能构造出海量练习题。它不是被教会“事实是什么”，而是被训练成“在这种上下文后，什么 token 最像会出现”。

模型最后输出的不是一个词，而是一整张候选 token 的分数表。通常会先得到一组 logits，可以粗略理解成“还没有归一化的倾向”。再用 softmax 把它们变成概率：

\[
P_i=\frac{\exp(z_i)}{\sum_j\exp(z_j)}
\]

这里 \(z_i\) 是第 \(i\) 个 token 的 logit，\(P_i\) 是它被选中的概率。softmax 有一个很直观的作用：把所有候选项压成一组总和为 1 的概率。高分 token 会更突出，低分 token 不会完全消失。

她用笔帽轻轻点了点分母。

“大家都在里面。只是有的声音比较大。”

Producer 问：“所以不是选一个绝对正确的词？”

“嗯。是先给所有可能性排队。然后，让其中一个走出来。”

这就解释了一个常见现象：模型既有稳定性，也有随机性。对于“法国的首都是”这种上下文，正确 token 的概率会非常集中；对于“请写一个有初星学园味道的开场”，合理续写可能有很多种，概率分布就会更分散。技术上这只是概率分布的形状；落到使用者眼里，就变成了“模型有时很确定，有时很有创意”。

Bengio 等人在神经概率语言模型中已经把“预测下一个词”与分布式词表示结合起来\cite{ch02bengio2003}。后来的 Transformer、GPT 和聊天助手都没有离开这条骨架；它们真正改变的，是上下文怎样表示、参数怎样训练、输出怎样被约束，以及模型怎样接触外部世界。

广把“预测”两个字又圈了一遍。

过去的软件常常需要人先选择功能：翻译按钮、摘要按钮、问答按钮。语言模型则把许多任务都改写成“接下来应该出现什么文本”。翻译是接出另一种语言，摘要是接出压缩版，代码解释是接出自然语言说明。

这也是为什么论文里常说它是 self-supervised learning。训练数据不需要人工逐条标注“这句话的类别是什么”。互联网规模的文本因此变成了巨大的练习册。它不是老师一题一题批改出来的，而是从人类留下的文字痕迹里自己构造题目。

当然，这种训练方式有代价。它学到的是“文本中经常怎样说”，不是“世界事实上怎样运转”。如果许多文本把某个错误说法重复很多遍，模型也可能学进去。如果训练数据里缺少某个小众领域，模型对它的把握就会薄。语言模型的数学目标很漂亮，但它和真实世界之间仍然隔着数据、语境和验证。

这就是它后来能成为通用接口的原因。

她在白板边写下一句话：

\begin{quote}
同一个训练目标，覆盖了很多任务的入口。
\end{quote}

很简单。

也很过分。

不过，把目标写成概率，只是给问题画出骨架。语言本身不会因此变简单。下一步要问的是：如果人想直接把语言规则写进机器里，为什么很快就会失败？

\section{语言太难：规则写不完，统计看不远}

大语言模型的“前世”并不是一开始就有神经网络。

早期聊天程序 ELIZA 是一个很好的入口。Weizenbaum 在 1966 年发表 ELIZA 论文时，展示了一个可以和人进行自然语言交流的程序\cite{ch02weizenbaum1966}。它最著名的 DOCTOR 脚本会模拟心理咨询式对话：用户说“我很难过”，程序可能把句子改写成“你为什么觉得自己很难过？”

这不是现代意义上的理解。它主要依赖模式匹配和替换规则。

所谓模式匹配，就是先在输入里找某种句式，再把句子的一部分填进预先写好的模板里。用户说“我觉得 X”，程序不需要知道 \(X\) 具体意味着什么，只要把它抽出来，套进“为什么觉得 X？”这样的回答里。规则系统的聪明，来自人事先写好的脚本；系统本身并不会从大量对话中更新自己的语言习惯。

Producer 问：“那 ELIZA 和现在的聊天模型，都是在装懂？”

“装懂这个词太快了。”广说，“ELIZA 是规则像懂。大语言模型是统计结构像懂。像的原因不同，危险也不同。”

\begin{figure}[htbp]
\centering
\begin{tikzpicture}[>=stealth, every node/.style={font=\small}]
  \node[draw, rounded corners, align=center, text width=2.6cm] (input) at (0,0) {用户输入\\我觉得 X};
  \node[draw, rounded corners, align=center, text width=2.8cm] (rule) at (4,0) {规则表\\匹配 ``我觉得''};
  \node[draw, rounded corners, align=center, text width=2.8cm] (output) at (8,0) {模板输出\\为什么觉得 X?};
  \draw[->] (input) -- (rule);
  \draw[->] (rule) -- (output);
\end{tikzpicture}
\caption{规则系统的直觉：先匹配句式，再套用模板。}
\end{figure}

它像一个背了很多反问句的人。接话很顺，但没有在心里形成世界模型。可是 ELIZA 的重要性正在这里：它很早就让人看到，只要语言表面足够像理解，人类就可能把理解投射上去。

这件事对今天仍然有提醒意义。很多人第一次和大语言模型聊天时，也会因为它的礼貌、连贯和自信而产生“它懂我”的感觉。ELIZA 时代靠规则制造这种错觉，今天靠大规模统计结构制造更强的错觉。机制不同，人的心理反应却有连续性。

规则系统的问题也很清楚。语言太柔软，场景太多，人又喜欢省略、反讽、打错字。规则写得少，覆盖不了；规则写得多，维护困难，还容易互相冲突。

例如“没关系”三个字，写进规则表时看起来很简单。可是在不同场景里，它可能表示真的原谅，也可能表示压抑不满，也可能只是礼貌结束话题。规则系统必须提前知道场景、语气、关系、前文。只靠句式匹配，很难判断这些看不见的东西。

后来统计语言模型开始接过接力棒。一个经典近似是 n-gram：

\[
P(w_t\mid w_{<t})\approx P(w_t\mid w_{t-n+1},\ldots,w_{t-1})
\]

意思是：预测当前词时，不看全部前文，只看前面 \(n-1\) 个词。这样问题变得可计算，也能从大语料里统计频率。n-gram 的 \(n\) 指的是连续 \(n\) 个词组成的片段；bi-gram 看两个词，tri-gram 看三个词。它不理解句子，只数“这种前缀后面，历史上常接什么”。

如果 \(n=2\)，模型只看前一个词；如果 \(n=3\)，就看前两个词。比如“语言 模型”后面常常接“可以”“训练”“生成”，模型就会给这些词更高概率。这个思路非常朴素，却把语言处理从手写规则推进到了数据驱动：我们不必把所有语法都写出来，可以让语料告诉机器哪些组合常见。

Producer 看着公式：“那只要把 \(n\) 变得很大，不就解决了？”

“看起来是这样。”广点头，“但是词表会膨胀，组合会爆炸，没见过的组合会越来越多。”

但短窗口有短窗口的命运。它看不远。

比如这句话：

\begin{quote}
那本放在图书馆最里面、封面已经褪色、被我拿来压住讲义的书，终于被 Producer 找到了。
\end{quote}

“找到”的对象是很早出现的“书”。如果模型只看前几个词，它很容易迷路。Bengio 等人指出，传统统计模型还会遇到维度灾难：词组合太多，很多合理句子在训练数据里根本没出现过\cite{ch02bengio2003}。

这叫数据稀疏。语言的组合太多了，哪怕语料很大，也不可能覆盖所有句子。一个人从没听过“我把训练记录夹在植物图鉴里”这句话，也能理解它；但如果统计模型完全依赖出现次数，这种新组合就会很难处理。它需要某种泛化能力：知道“训练记录”和“笔记”相近，“夹在书里”和“放进书页”相近，即使原句没出现过，也能推断它的合理性。

于是问题变成了：能不能不要把每个词都当作孤立符号？

能不能让词语之间拥有可学习的距离？

\subsection*{词不是孤岛：从统计频率到连续表示}

词向量的想法，是把词从离散编号变成连续空间里的点。

“猫”“狗”“动物”会因为上下文相似而靠近；“国王”“王后”“男人”“女人”之间可能形成某种方向关系。Mikolov 等人的 word2vec 工作让这种直觉变得非常有影响力：它用高效架构从大规模语料中学习连续词表示，并在词相似度与语法语义关系上取得很强效果\cite{ch02mikolov2013}。

这背后有一个很重要的语言学直觉：词的意义可以从它出现的邻居里看出来。一个词如果经常和“训练”“舞台”“出汗”“节拍”一起出现，它就会被拉向偶像活动和身体练习的语境；一个词如果经常和“参数”“损失”“梯度”一起出现，它就会被拉向机器学习的语境。模型没有读词典，却从上下文里学到词的用法。

在神经网络里，词向量通常来自一个嵌入矩阵。词表里每个 token 对应一行向量。输入 token 之后，模型先查表得到向量，再把向量送入后续网络。这样，“猫”和“狗”不再只是两个互不相干的编号，而是两个可以比较距离、相似度和方向的点。

Producer 问：“所以向量就是给词换了一个编号？”

“不是。”广说，“编号只负责区分。向量还负责靠近和远离。”

如果“猫”和“狗”只是编号 1357 和 2468，模型不知道它们有什么关系。词向量把它们放进连续空间里，相似词会因为上下文相似而靠近。这样，模型遇到训练集中没见过的新组合时，也可以借助相近词的经验做一点泛化。它不再只背“出现过的词串”，而是开始学习词和词之间的可计算关系。

衡量两个词向量相似，常用余弦相似度：

\[
\cos(v_a,v_b)=\frac{v_a\cdot v_b}{\|v_a\|\|v_b\|}
\]

如果两个向量方向接近，余弦值就高。你可以把它想成：两个词是否经常出现在相似语境里。

\begin{figure}[htbp]
\centering
\begin{tikzpicture}[>=stealth, every node/.style={font=\small}]
  \draw[->] (-0.2,0) -- (7.2,0) node[right] {维度 1};
  \draw[->] (0,-0.2) -- (0,3.6) node[above] {维度 2};
  \node[circle, fill=accent, inner sep=2pt, label=above:{猫}] at (1.4,2.4) {};
  \node[circle, fill=accent, inner sep=2pt, label=right:{狗}] at (2.0,2.1) {};
  \node[circle, fill=accent, inner sep=2pt, label=above:{动物}] at (1.7,3.0) {};
  \node[circle, fill=primary, inner sep=2pt, label=above:{代码}] at (5.2,1.1) {};
  \node[circle, fill=primary, inner sep=2pt, label=right:{函数}] at (5.8,1.5) {};
  \node[circle, fill=primary, inner sep=2pt, label=below:{编译}] at (4.8,0.6) {};
\end{tikzpicture}
\caption{词向量的直觉示意：语境相似的词，在空间中更接近。}
\end{figure}

\section{模型长出身体：注意力、预训练与规模}

词向量解决了“词不是孤岛”的问题，但文本还有顺序。句子不是一袋词，而是一条时间线。

“Producer 批评了我”和“我批评了 Producer”用到几乎同一批词，但意思完全不同。词向量能表示词的关系，却不能单独解决顺序和句法。模型还需要知道谁在前，谁在后，哪个修饰哪个，哪个动作指向哪个对象。

循环神经网络尝试按顺序读文本。一个非常简化的状态更新可以写成：

\[
h_t=f(Wx_t+Uh_{t-1})
\]

这里 \(h_t\) 是读到第 \(t\) 个位置时的隐藏状态。它既看当前输入 \(x_t\)，也看上一步的记忆 \(h_{t-1}\)。LSTM 则进一步引入门控机制，试图缓解普通循环网络在长序列里容易遗忘的问题\cite{ch02lstm1997}。

LSTM 的“门”可以粗略理解成三个决定：哪些旧信息要留下，哪些新信息要写入，哪些状态要输出。它比普通 RNN 更像一个会整理笔记的人，不是每读一个词就把所有信息混在一起，而是尝试决定什么值得记住。

这一段可以很短，因为 RNN 不是终点。它的贡献，是把“顺序读”和“带着记忆读”变成了神经网络里的自然操作；它的限制，也正来自这里。它必须一步一步读，不能轻易跳过时间线。

这像边读边做笔记。问题是，笔记本越传越厚，训练越慢；句子越长，早期信息越难稳定保留。模型想记住前文，但前文隔得太远时，它还是会累。

更麻烦的是，循环结构天然按时间一步一步计算。第 \(t\) 步必须等第 \(t-1\) 步算完。对于短句，这还可以接受；对于互联网规模的文本训练，这会成为巨大瓶颈。模型不只是要会记忆，还要能被成千上万张计算卡高效训练。

广说：

\begin{quote}
会累不是失败。\\
是结构在告诉我们，下一步该换方法了。
\end{quote}

说完以后，广发现它很像训练记录，也很像模型史。

\subsection*{注意力成为骨架：让每个位置看向别处}

2017 年，Vaswani 等人发表 \textit{Attention Is All You Need}，提出 Transformer 架构\cite{ch02vaswani2017}。这篇论文的关键主张很大胆：处理序列不一定要依赖循环结构，也不一定要依赖卷积；注意力机制本身就可以成为核心骨架。

注意力的基本公式是：

\[
\mathrm{Attention}(Q,K,V)=
\mathrm{softmax}\left(\frac{QK^T}{\sqrt{d_k}}\right)V
\]

公式里有三个角色：Query、Key、Value。可以把 Query 理解成“我现在在找什么”，Key 是“每个位置提供的索引”，Value 是“真正要取出的信息”。模型用 \(QK^T\) 计算当前位置和其他位置的相关程度，再用 softmax 变成权重，最后加权汇总 Value。

在实现上，\(Q,K,V\) 通常都来自同一批输入向量，只是经过不同线性变换：

\[
Q=XW_Q,\quad K=XW_K,\quad V=XW_V
\]

这里没有人手工规定“主语应该看谓语”或“代词应该看名词”。模型要在训练中自己学会：哪些位置之间互相有用。不同层、不同注意力头会形成不同观察方式：有的头关注相邻词，有的头关注标点，有的头追踪长距离指代。

Producer 把“注意力”圈起来：“这是不是就等于理解？”

“不等于。”广说，“注意力只是计算关联的一种方法。它让位置之间更容易交换信息，不保证交换出来的东西一定正确。”

self-attention 的重点在 self：同一段序列里的 token 彼此计算相关性。每个位置都可以形成自己的 Query，去和其他位置的 Key 比较，再按权重汇总 Value。这样，模型不必像 RNN 那样把前文压进一条递推状态，而是可以在一层计算里让许多位置直接建立联系。

这听起来抽象。换成课堂比喻就简单一点。

老师问：“这句话里的‘它’指什么？”

学生不能只盯着“它”这个字。需要回头看前文：是“语言模型”，还是“训练数据”，还是“注意力机制”？注意力就是让模型在不同位置之间分配目光。

\begin{figure}[htbp]
\centering
\begin{tikzpicture}[every node/.style={font=\small}]
  \node at (0,2.6) {查询词：它};
  \foreach \x/\word/\shade in {0/模型/30,1.6/读过/10,3.2/很多/15,4.8/文本/45,6.4/所以/8,8.0/它/80} {
    \node[draw, fill=accent!\shade, minimum width=1.1cm, minimum height=0.65cm] at (\x,1.4) {\word};
  }
  \node[align=center, text width=7.8cm] at (4,0.25) {颜色越深，表示当前 token 对该位置分配的注意力越高。};
\end{tikzpicture}
\caption{注意力热力图的简化示意。真实模型会在多层、多头中学习更复杂的关系。}
\end{figure}

Transformer 还需要位置编码。因为它不按 RNN 那样一步一步读，所以必须把位置信息注入 token 表示里。原论文使用正弦、余弦位置编码，让模型知道“这个 token 在第几个位置”\cite{ch02vaswani2017}。

它适合大规模训练，原因有三个。

第一，注意力可以直接建立远距离联系。第二，它比循环结构更容易并行计算。第三，多层堆叠后，模型可以逐层组合局部和全局信息。

这里还要加一个边界。Transformer 不是永远让所有 token 随便互看。BERT 这类编码器模型通常可以同时看左右文；GPT 这类自回归生成模型会使用 causal mask，让当前位置不能看到未来 token。也就是说，它能在允许的范围内直接看向别的位置，但生成时仍然要遵守“从左到右”的时间方向。

“多头注意力”也值得多说一句。单个注意力头像一束视线，多头注意力则像从多个角度同时看同一句话。一个头看语法关系，一个头看主题词，一个头看引用对象，一个头看格式边界。它们的结果拼接起来，再经过前馈网络处理。这样，模型不是用一种目光理解文本，而是用许多目光同时扫描。

广把四支不同颜色的笔并排放在桌上。

“同一句话，可以有很多种看法。”

她看向 Producer。

“你看我走路，会看步幅。看呼吸。看脸色。看我是不是在说谎。”

“你经常在说‘没事’的时候说谎。”

“嗯。所以需要多头注意力。”

这也是 Transformer 和早期模型气质不同的地方。RNN 像一个人沿着走廊一步步走，边走边记；Transformer 像把整段话摊在桌上，让每个词在规则允许的范围内直接看向其他词。对于长文本和大规模训练，这种结构上的变化非常关键。

这就是现代大语言模型真正的转折点。

不是因为它第一次会说话，而是因为它终于拥有了一副可以被放大的身体。

\subsection*{预训练：先读很多书}

Transformer 之后，预训练成为主线。

BERT 代表一种理解型路线。它使用双向 Transformer，从左右两侧上下文中学习语言表示。论文强调，BERT 可以在无标注文本上预训练，然后只加很少的任务层，就在问答、自然语言推理等多种任务上取得强效果\cite{ch02bert2018}。

可以把 BERT 想成做完形填空的学生：

\begin{quote}
白板上写着 [MASK]。
\end{quote}

模型要根据左右文判断，遮住的词可能是什么。它擅长理解句子内部关系。

BERT 的代表性任务叫 masked language modeling：把输入里一部分 token 遮住，让模型根据左右文恢复它。因为它做题时能同时看前后，所以很适合给整段文本做判断、匹配、分类和抽取。

Producer 问：“那它不能写文章吗？”

“不是完全不能。”广说，“只是它出生时练的题型，不是从左到右续写。”

GPT 路线更偏生成。GPT-2 和 GPT-3 使用自回归语言建模：只看左边上下文，从左到右生成后续 token\cite{ch02gpt22019,ch02gpt32020}。自回归的意思是，模型把自己刚生成的 token 也放回上下文里，再预测下一个。它像一边走一边铺轨道：前一步写出来的东西，会变成下一步的条件。

这两条路线的差别，不是好坏之分，而是训练目标不同。BERT 在预训练时能同时看左边和右边，所以特别适合做分类、匹配、抽取式问答等“读懂整段文本再判断”的任务。GPT 在生成时只能看左边，所以天然适合继续写：它每一步都只需要回答“接下来是什么”。

后来聊天式大模型大多沿着 GPT 式生成路线发展。原因很实际：对话、写作、代码、计划，最后都要生成一串新 token。理解当然重要，但它被包进生成过程里。模型必须先从上下文里读出任务，再继续写出合适输出。

\begin{figure}[htbp]
\centering
\begin{tikzpicture}[>=stealth, every node/.style={font=\small}]
  \node[draw, rounded corners, align=center, text width=3.2cm] (bert) at (0,0) {BERT\\看左右文\\填空理解};
  \node[draw, rounded corners, align=center, text width=3.2cm] (gpt) at (5.2,0) {GPT\\看左文\\继续生成};
  \node[align=center, text width=3.3cm] at (0,-1.4) {白板上 [MASK]};
  \node[align=center, text width=3.3cm] at (5.2,-1.4) {白板上 → 写下};
\end{tikzpicture}
\caption{BERT 与 GPT 的直觉差异：一个偏填空理解，一个偏顺序生成。}
\end{figure}

预训练的意义，是让模型先在大量文本中学到通用语言规律，再迁移到具体任务。过去很多 NLP 系统要为每个任务单独设计结构；预训练模型则把大量知识和语言模式压进参数里。

这里的“知识”要打引号。它不是数据库里的整齐条目，而是从文本分布中学到的关联。模型知道很多事情，是因为相关表述在训练数据中反复出现；它也会混淆很多事情，因为文本本身有噪声、矛盾和过时信息。

预训练还有一个重要后果：模型内部没有清晰的“知识出处”。人类读书时可以说“我在某本书第几章看到过”，数据库可以返回某条记录，预训练模型却只是把大量文本压进参数。它回答问题时，不是在翻开某一页，而是在参数空间里生成一个看起来合理的延续。这一点会在后面的 RAG 部分变得很重要。

广把预训练称作“漫长的旁听”。

没有老师逐题讲解。

只是坐在那里，听了非常多句子。

嗯。

很像我旁听体能课。

脑子懂了。

身体还没有。

“听得太久，也会累吧？”Producer 问。

“会。”广说，“但是累不一定不好。说明容量正在被使用。”

\subsection*{上下文里学任务：规模带来的临场适应}

GPT-3 的重要性，不只在于参数大。它更醒目的地方，是把 few-shot learning 变成了大语言模型时代的标志性能力\cite{ch02gpt32020}。

所谓 few-shot，就是不用重新训练模型，只在提示词里给几个例子。比如：

\begin{quote}
中文：早上好。\\
日文：おはよう。\\
中文：谢谢。\\
日文：ありがとう。\\
中文：今天的训练很难。\\
日文：
\end{quote}

模型看到这种格式，就会继续补出对应答案。它没有为“这个翻译任务”单独更新参数，而是在上下文里识别任务模式。

这件事后来被称作 in-context learning。它像一种临场适应：参数不变，但上下文改变了模型当前的行为。你给它三个翻译例子，它就进入翻译模式；你给它几条 JSON 样例，它就尝试继续生成 JSON；你给它一段角色设定，它就调整语气和词汇。

这也是提示词为什么重要。提示词不是对模型“下命令”这么简单，它更像临时搭了一个小环境。环境里有规则、示例、格式、限制和语气。模型根据这个环境继续生成。环境越清楚，模型越容易落到正确轨道；环境越含糊，它就越可能走偏。

Producer 问：“所以 prompt engineering 是给模型写说明书？”

“不只是说明书。”广说，“更像布置考场。”

说明书告诉它要做什么，例子告诉它题型是什么，格式告诉它答案长什么样，角色和语气告诉它要站在哪里说话。强模型会从这些上下文线索里推断当前任务。弱模型也会照着格式走，但可能只学到表面。

\begin{figure}[htbp]
\centering
\begin{tikzpicture}[>=stealth, every node/.style={font=\small}]
  \node[draw, rounded corners, align=center, text width=2.1cm] (ex1) at (0,0) {示例 1\\输入→输出};
  \node[draw, rounded corners, align=center, text width=2.1cm] (ex2) at (2.8,0) {示例 2\\输入→输出};
  \node[draw, rounded corners, align=center, text width=2.1cm] (q) at (5.6,0) {新问题\\输入→?};
  \node[draw, rounded corners, align=center, text width=2.1cm] (a) at (8.4,0) {模型续写\\输出};
  \draw[->] (ex1) -- (ex2);
  \draw[->] (ex2) -- (q);
  \draw[->] (q) -- (a);
\end{tikzpicture}
\caption{少样本学习的直觉：任务说明和例子直接放进上下文。}
\end{figure}

这件事很像临时看题型。老师没有重新训练学生的大脑，只是在卷面上给了两道样题。足够强的学生会看出格式，然后照着做第三题。

GPT-3 论文同时也很谨慎。它指出模型在一些任务上仍然挣扎，也讨论了大规模网页语料带来的方法问题和社会影响\cite{ch02gpt32020}。这很重要：大，不等于稳；会举一反三，不等于每次都对。

few-shot 的能力尤其容易让人高估模型。因为它看起来像“读懂了新任务”。但很多时候，它是在捕捉格式、词面线索和相似任务模式。简单任务上这足够好，复杂任务里就可能露出问题：题目条件稍微变形，示例带有误导，输出格式和事实正确性同时受约束，模型就可能只学到了表面形式。

所以，这一节的结论不是“模型突然会思考了”。

更准确的说法是：当模型足够大、训练足够广，它开始能把上下文本身当作任务说明。

广说：“它不是读懂了考试制度。”

“它只是很会看题面。”

这里还欠一个问题：为什么“足够大、训练足够广”会带来这种临场适应？它不是魔法，而是规模、数据和计算共同把损失压低以后，许多语言模式开始在同一个模型里相互借力。

\subsection*{规模律：大，不只是参数大}

为什么大模型会变强？业界论文里有一条重要线索叫 scaling laws。

Kaplan 等人的工作发现，语言模型损失会随着模型规模、数据规模和训练计算量呈现较规律的幂律下降趋势\cite{ch02scaling2020}。简化地说，可以写成一种直觉形式：

\[
L(N,D,C)\downarrow
\quad\mathrm{when}\quad
N,D,C\uparrow
\]

其中 \(N\) 是参数量，\(D\) 是训练数据量，\(C\) 是计算量。实际论文里的拟合更细，但科普时抓住这个方向就够：变大通常会带来更低损失、更强能力。

这里的“损失下降”不是一个抽象数字。它意味着模型在预测下一个 token 时更少惊讶。很多能力就是从这种“不那么惊讶”里浮出来的：它更熟悉代码结构，更熟悉论文写法，更熟悉多轮对话，更熟悉人类如何一步步解释问题。

可是，如果只记住“越大越好”，就会误读这条线。

Chinchilla 论文进一步指出，在固定计算预算下，很多大语言模型其实是“训练不足”的：参数很多，但数据不够。它建议模型参数量和训练 token 数应当更均衡地增长，并展示了 70B 参数的 Chinchilla 在相同计算预算下击败更大模型的结果\cite{ch02chinchilla2022}。

\begin{figure}[htbp]
\centering
\begin{tikzpicture}[>=stealth, every node/.style={font=\small}]
  \draw[->] (0,0) -- (7.4,0) node[right] {参数量 \(N\)};
  \draw[->] (0,0) -- (0,3.6) node[above] {训练数据 \(D\)};
  \draw[thick, color=accent] (0.7,0.6) -- (6.4,3.1);
  \node[align=center, text width=2.5cm] at (4.9,2.1) {较均衡\\训练更充分};
  \node[draw, rounded corners, align=center, text width=2.3cm] at (5.2,0.8) {参数大\\数据不足};
  \node[draw, rounded corners, align=center, text width=2.3cm] at (1.6,2.8) {数据多\\模型太小};
\end{tikzpicture}
\caption{Chinchilla 式直觉：固定预算下，模型大小和训练数据要匹配。}
\end{figure}

LLaMA 的思路也延续了这种更重视训练效率的方向。它强调用更多 token 训练较小模型，展示了开放权重基础模型可以在较低推理成本下取得很强表现\cite{ch02llama2023}。后来 Llama 3.1 等模型进一步把“开放权重模型也能接近前沿能力”变成产业事实之一\cite{ch02llama312024}。这里说开放权重，而不是简单说开源，是因为模型权重、训练数据、训练代码和使用许可的开放程度并不总是相同。

这对产业也很关键。一个最大、最贵的模型可能在榜单上最好，但真实应用还要考虑延迟、成本、隐私、部署环境和可控性。客服系统、手机输入法、企业知识库、嵌入式设备，不一定都需要最大模型。它们需要的是在目标任务上足够好、足够快、足够便宜、足够安全。

Producer 问：“那模型越小越好吗？”

“也不是。”广说，“只是问题变成了匹配。”

大模型像总科室，见多识广，但贵、慢、部署难；小模型像专门训练过的值日生，在固定场景里可能更快、更稳、更便宜。规模律讲的是趋势，不是替你做产品决策。真正的问题是：任务需要多少能力，数据能不能支撑，预算和延迟能不能接受。

因此，规模不是单一旋钮。它像训练菜单。

只加重量，不加恢复，会受伤。只加练习量，不提升动作质量，也会卡住。模型也一样：参数、数据、算力、训练配方、评测方式，都在共同决定它最后的能力。

Producer 在旁边把“休息”两个字写得很大。

广看了一眼。

“我知道。”

“你每次都说知道。”

“因为每次都还没学会。嗯，很适合复习。”

嗯。

不是只要变大就好。

要正确地变大。

可是，正确地变大之后，模型仍然只是在续写。它可以看懂题面，却未必知道自己应该认真答题；它可以接出一段像答案的文字，却未必站在用户这一边完成任务。下一道墙，就在这里。

\section{从续写到助手：对齐、幻觉与检索}

原始预训练模型会续写文本，但它不一定会“听话”。

你问它“请解释这段代码”，它可能解释，也可能续写成一篇博客，也可能编造上下文。原因很简单：预训练目标只要求它预测文本分布，没有明确告诉它“用户希望你完成任务”。

InstructGPT 论文把这个问题讲得很清楚：更大的语言模型并不天然更符合用户意图，可能不真实、有害，或者根本不帮忙\cite{ch02instructgpt2022}。他们使用人类反馈进行训练，大致分三步：

\begin{enumerate}
  \item 收集人工示范，用监督学习微调模型。
  \item 收集人类对多个回答的排序，训练奖励模型。
  \item 用强化学习优化模型，让它生成更受人类偏好的回答。
\end{enumerate}

\begin{figure}[htbp]
\centering
\begin{tikzpicture}[>=stealth, every node/.style={font=\small}]
  \node[draw, rounded corners, align=center, text width=2.3cm] (sft) at (0,0) {人工示范\\SFT};
  \node[draw, rounded corners, align=center, text width=2.5cm] (rm) at (3.4,0) {答案排序\\奖励模型};
  \node[draw, rounded corners, align=center, text width=2.5cm] (rl) at (6.9,0) {强化学习\\优化策略};
  \draw[->] (sft) -- (rm);
  \draw[->] (rm) -- (rl);
\end{tikzpicture}
\caption{RLHF 的三阶段直觉：先学示范，再学偏好，最后按偏好优化。}
\end{figure}

一个直觉式的 RLHF 目标可以写成：

\[
\max_\pi \; r_\phi(x,y)-\beta D_{\mathrm{KL}}(\pi(y\mid x)\,\|\,\pi_0(y\mid x))
\]

这里 \(r_\phi\) 表示奖励模型给回答的分数，后面的 KL 项约束新模型不要离原模型太远。实际训练细节会更复杂；这个式子只是抓住一个核心直觉：让模型更接近人类偏好的回答，同时不要偏离原本语言能力太多。

直觉上，它像 Producer 给广安排训练：可以把她推向更好的表现，但不能把人推到完全失控。

“准确的严厉，是训练。”广说。

“没有边界的严厉呢？”

“那只是噪声。没有学习价值。”

这里最容易误解的是“奖励”。奖励模型不是事实裁判，也不是道德本身。它只是从人类偏好数据中学到“哪种回答更像人类标注者喜欢的回答”。如果标注标准好，它会帮助模型更有用、更诚实、更安全；如果标准含糊，模型也可能学会讨好、回避或说得更漂亮。

Producer 问：“那 RLHF 会让模型变真实吗？”

“会改善。”广说，“但不能保证。”

人类偏好可以奖励“承认不知道”，也可以奖励“说话有条理”。可是事实正确性还需要数据、检索、工具、评测和产品约束一起工作。一个回答听起来诚恳，不等于它引用的论文真的存在。对齐让模型更像助手，不让模型自动变成事实数据库。

所以 RLHF 的意义不是让模型获得真理，而是把“续写互联网文本”的模型拉向“回答用户问题”的模型。它把原本开放的文本分布，收束到更符合助手角色的区域：更愿意遵循指令，更少输出危险内容，更常承认不确定。

LaMDA 论文从对话模型角度讨论了质量、安全与事实扎根等指标\cite{ch02lamda2022}。ChatGPT 的发布则把这种指令跟随与对话交互推向大众产品\cite{ch02chatgpt2022}。

对普通使用者来说，这一步比参数量更有体感。一个没有对齐的强模型像非常会写但不听题的学生；一个经过指令微调和人类反馈的模型，才像愿意按照题目作答的助手。ChatGPT 的“产品震感”很大程度来自这里：人第一次可以不写代码、不设计任务头，只用自然语言描述目标，就驱动模型工作。

所以，ChatGPT 不是“语言模型突然会聊天”。它更像一个预训练模型被指令微调、人类反馈和产品化对话界面慢慢推到桌前，终于开始认真听题。

广在白板角落补了一行：

\begin{quote}
它不是变成人。\\
它只是终于知道，题目要求不是“继续写”，而是“帮我做”。
\end{quote}

\subsection*{会帮忙，不等于说真话}

大语言模型最危险的地方，不是它什么都不会。

而是它经常说得很像会。

模型训练目标是生成高概率文本，不是保证事实正确。它可能把真实术语、相似人名、常见论文格式和不存在的细节拼在一起，生成一个看起来很可靠的错误答案。这就是常说的幻觉。

幻觉的麻烦在于，它不是总以“明显错误”的样子出现。低质量答案很容易发现；高质量幻觉反而危险。它可能格式正确，语气稳重，引用看起来像真的，推理过程也很顺。如果只看表达流畅度，很容易被它带走。

解决幻觉不能只靠“让模型更诚实”。很多时候，模型需要外部资料。

RAG，也就是检索增强生成，把语言模型和信息检索结合起来。Lewis 等人的 RAG 工作把参数化生成模型与非参数化记忆结合，在知识密集型任务上取得改进\cite{ch02rag2020}。广在旁边写了更短的说法：不要只凭记忆答题，先查资料，再带着资料回答。

RAG 的价值有两层。第一，它能补充新知识。模型训练完成后，参数里的知识会慢慢过时，而检索系统可以访问更新的文档。第二，它能提供可检查的证据。模型不只是给一个答案，还可以把相关资料片段带出来，让人知道它依据什么回答。

Producer 问：“所以接了 RAG，就不会胡编了？”

“不会。”广说，“只是多了一条能被检查的路。”

RAG 至少有三件事要同时成立。第一，检索器要找对资料；第二，资料片段要真的进入模型上下文；第三，生成答案时要忠实使用这些资料。任何一层出错，最后都可能变成“带着资料胡编”。有时模型找到了正确文档，却引用错段落；有时资料本身过期；有时多个来源互相冲突，模型选择了看起来更顺口的那个。

\begin{figure}[htbp]
\centering
\begin{tikzpicture}[>=stealth, every node/.style={font=\small}]
  \node[draw, rounded corners, align=center, text width=2cm] (q) at (0,0) {问题};
  \node[draw, rounded corners, align=center, text width=2.4cm] (ret) at (2.8,0) {检索器\\找资料};
  \node[draw, rounded corners, align=center, text width=2.5cm] (ctx) at (5.9,0) {上下文\\证据片段};
  \node[draw, rounded corners, align=center, text width=2.2cm] (gen) at (9,0) {模型\\生成回答};
  \draw[->] (q) -- (ret);
  \draw[->] (ret) -- (ctx);
  \draw[->] (ctx) -- (gen);
\end{tikzpicture}
\caption{RAG 的直觉：先找证据，再生成回答。}
\end{figure}

GPT-4 技术报告也反复强调能力评测、安全评估与局限性\cite{ch02gpt42023}。Constitutional AI 则探索用一组原则和 AI 反馈来改进无害性，减少完全依赖人类逐条标注的成本\cite{ch02constitutional2022}。

但 RAG 也不是万能药。检索器可能找错资料，资料本身可能过时，拼进上下文的证据可能被模型误读。更现实的问题是：当证据很多、互相冲突、质量不一时，模型仍然需要判断。可靠系统不能只写“接入 RAG”四个字，还要设计文档清洗、权限控制、引用格式、冲突处理、答案溯源和人工复核。

这一段的重点是：可靠性不是一个按钮。

它需要多层机制：更好的训练数据，更好的对齐，更好的检索，更明确的引用，更严格的评测，以及人在关键场景里的复核。

广在这里停顿了一下。

“它很强。所以更要确认。”

这句话不是保守。

是使用强工具时最基本的礼貌。

她把讲义翻到引用列表那一页。

“资料也会说谎。或者过期。或者只说了一半。”

Producer 点头：“所以要看来源。”

“嗯。プロデューサー也是。”

“我？”

“你说‘今天只练一点’，也需要验证。”

到这里，模型已经从“会接话”走到“会回答”，又被事实和来源拉住。但现实工作不只是问答。很多任务需要查资料、算结果、改文件、操作系统界面。于是，语言模型开始从答案走向行动。

\section{从回答到行动：工具、Agent 与边界}

当模型能稳定对话后，人们很快发现：只让它直接给答案还不够。很多任务需要中间步骤。

Chain-of-Thought 论文提出，在提示中给出推理过程示例，可以显著提升大模型在算术、符号和常识推理任务上的表现\cite{ch02cot2022}。它的直觉很简单：不要只让模型写最终答案，而是让它先把问题拆开。

这里的 chain 不是神秘的大脑链条，而是可见的中间文本。模型先生成“我应该怎样分解问题”，再用这些文字继续生成后面的步骤。对自回归模型来说，前面生成的每一句话都会成为后面预测的上下文。因此，中间步骤写得好，后面更容易落到正确答案；中间步骤写错，也会把错误传下去。

例如：

\begin{quote}
问题：练习室有 3 个水杯，我又拿来 2 个，Producer 收走 1 个，还剩几个？\\
推理：先有 3 个，加 2 个变 5 个，收走 1 个，还剩 4 个。\\
答案：4。
\end{quote}

注意，这不意味着模型真的像人一样在心里思考。它只是生成了中间文本，而这些中间文本会帮助后续 token 更容易落到正确轨道上。

因此，链式思维更像一种外化的草稿纸。对复杂问题，人类也常常需要把步骤写出来。模型生成中间步骤后，后续生成可以引用前面的中间结果，减少一步跳到答案时的压力。不过，草稿纸也可能写错。一步错，后面每一步都可能跟着错。

ReAct 进一步把 reasoning 和 acting 交替起来：模型一边写推理痕迹，一边执行动作，从外部环境获得观察结果\cite{ch02react2022}。Toolformer 则探索让模型学习什么时候调用外部 API、传什么参数、如何把工具返回结果融入后续预测\cite{ch02toolformer2023}。

Producer 问：“只要会调用工具，就是 agent 吗？”

“还不够。”广说，“会拿工具，和会完成任务，中间差很多。”

从这里开始，模型不再只是语言空间里的续写器。它可以把“需要知道今天的天气”转成查询工具，把“要分析这份表格”转成代码执行，把“要安排会议”转成日历操作。语言变成控制层，工具变成手脚。

\begin{figure}[htbp]
\centering
\begin{tikzpicture}[>=stealth, every node/.style={font=\small}]
  \node[draw, rounded corners, align=center, text width=2.3cm] (thought) at (0,0) {Thought\\需要查什么};
  \node[draw, rounded corners, align=center, text width=2.3cm] (action) at (3.5,0) {Action\\调用工具};
  \node[draw, rounded corners, align=center, text width=2.3cm] (obs) at (7,0) {Observation\\得到结果};
  \draw[->] (thought) -- (action);
  \draw[->] (action) -- (obs);
  \draw[->] (obs) to[bend left=35] (thought);
\end{tikzpicture}
\caption{ReAct 式循环：推理、行动、观察交替进行。}
\end{figure}

这就是 AI agent 的雏形：模型不只是回答，还能规划、查资料、调用工具、观察结果、继续修正。

agent 的关键不是“自主”两个字，而是闭环。只生成计划，不观察执行结果，还是纸上谈兵；只调用工具，不会解释目标，也只是自动脚本。真正有用的 agent 要能把目标拆成步骤，在每步之后读取环境反馈，再决定下一步。

当然，这也放大了风险。一个只会写错答案的模型，最多误导你；一个会调用工具的模型，可能发错邮件、删错文件、查错数据。所以 agent 的核心不是“让模型自由行动”，而是给它边界、权限、回滚和审计。

到 2025 年以后，这条路线已经进入产品层。OpenAI 的 Computer-Using Agent 把视觉输入、鼠标键盘操作和网页任务结合起来，说明 agent 不只是调用一个天气 API，也可能直接操作图形界面\cite{ch02cua2025}。这让边界问题更严肃：模型不只是在文本里犯错，而可能在真实软件里点错按钮。

广说：

\begin{quote}
能动起来，不代表可以乱跑。
\end{quote}

プロデューサー 会先确认身体状态。

系统也一样，要先确认安全边界。

\subsection*{已经显形的未来：预测语言碰到世界边缘}

讲义写到这里时，大语言模型的“未来”已经不完全是未来。几条路线正在同时产品化。

先是多模态。模型不只处理文字，也处理图像、音频、视频和结构化数据。GPT-4 技术报告已经展示了更强的综合能力，虽然报告没有披露完整训练细节\cite{ch02gpt42023}。后续产品继续把文字、图像、语音、屏幕和文件放进同一个交互入口。

多模态的意义，不只是“模型能看图”。更重要的是，语言可以成为不同信息形态之间的中转层。你可以指着一张图问“这里为什么不对”，也可以让模型把视频里的步骤整理成文字，把表格里的异常解释成自然语言。语言模型开始从文字世界伸向感知世界。

Producer 问：“那它是不是开始看见世界了？”

“看见了一些输入。”广说，“不是拥有世界。”

图像、音频和屏幕只是新的上下文形态。模型能处理它们，不代表它自动理解现实因果、物理限制和社会责任。多模态让模型接近世界，但世界不是只靠输入格式就能被完全装进上下文窗口。

另一个方向是开放权重和本地化。LLaMA 等模型说明，强模型不一定只存在于少数封闭系统里；更小、更高效的模型可以部署在企业内部、个人电脑或移动设备上\cite{ch02llama2023,ch02llama312024}。

本地化会改变很多应用的边界。医疗、法律、企业内部知识库、个人笔记，都有隐私和合规要求。不是所有数据都适合发到远程服务。小模型、领域模型和本地推理会成为大模型生态的重要部分。它们未必在所有通用榜单上最好，却可能在固定流程里更可控。

还有推理强化。DeepSeek-R1 展示了通过大规模强化学习激励推理能力的路线，尤其强调在没有先做监督微调的情况下，R1-Zero 也能出现显著推理能力\cite{ch02deepseekr12025}。OpenAI 的 o3 和 o4-mini 等推理模型也把“先分解、再求解、必要时使用工具”的产品形态推到更显眼的位置\cite{ch02openai2025o3}。

这条路线让“训练模型回答什么”变成“训练模型怎样到达答案”。如果奖励只看最终对错，模型可能发展出更长的试探、验证和修正过程。数学、编程、逻辑题尤其适合这类训练，因为答案可以相对明确地检查。

但推理模型也容易制造新的误会。长答案不一定更可靠，长推理不一定等于真推理。真正要看的不是它写了多少步骤，而是步骤能不能验证，工具结果是否真实，结论是否和证据一致。

最后是工具化和工作流化。模型会越来越少以单独聊天窗口的形态出现，越来越多嵌入编辑器、表格、设计工具、客服系统、编程环境和机器人流程中。到 GPT-5.5 这样的节点，模型能力已经和长上下文、工具调用、代码执行、文件处理、协作界面更紧密地绑在一起\cite{ch02gpt552026}。

这意味着未来的大语言模型未必总是以“一个人格化助手”的形式出现。很多时候，它会安静地藏在功能后面：自动补全一段说明，检查一条 SQL，给一封邮件换语气，把杂乱记录变成结构化字段。它越成熟，越可能变得不显眼。

但有一个问题不会消失：

\begin{quote}
当机器越来越会使用语言时，人类要怎样判断它是真的可靠，还是只是说得像可靠？
\end{quote}

这不是纯技术问题。它涉及教育、劳动、创作、法律、隐私、版权和责任分配。

所以，这堂课的结论要克制一点。

大语言模型不是魔法。它是一个把“预测下一个 token”训练到极限附近，再叠加数据、架构、反馈、检索、工具和产品约束的系统。因为语言连接了知识、工作、关系和世界，所以预测语言，也就碰到了世界的边缘。

广合上讲义，动作很轻，像是怕惊动纸页。

“边缘很危险。”

她想了想，又补上一句：

“也很适合上课。”

Producer 从她手边收走马克笔。

“今天到这里。”

“还可以再讲一点。”

“不可以。边界。”

广眨了眨眼。

“嗯。被正确拒绝了。”

她把这句话写进页边。

“下次继续。”

\begin{thebibliography}{99}
\bibitem{ch02weizenbaum1966} J. Weizenbaum, ``ELIZA--A computer program for the study of natural language communication between man and machine,'' \textit{Communications of the ACM}, vol. 9, no. 1, pp. 36--45, 1966. [Online]. Available: \href{https://doi.org/10.1145/365153.365168}{doi.org/10.1145/365153.365168}
\bibitem{ch02bengio2003} Y. Bengio, R. Ducharme, P. Vincent, and C. Jauvin, ``A neural probabilistic language model,'' \textit{Journal of Machine Learning Research}, vol. 3, pp. 1137--1155, 2003. [Online]. Available: \href{https://jmlr.org/papers/v3/bengio03a.html}{jmlr.org/papers/v3/bengio03a.html}
\bibitem{ch02lstm1997} S. Hochreiter and J. Schmidhuber, ``Long short-term memory,'' \textit{Neural Computation}, vol. 9, no. 8, pp. 1735--1780, 1997. [Online]. Available: \href{https://doi.org/10.1162/neco.1997.9.8.1735}{doi.org/10.1162/neco.1997.9.8.1735}
\bibitem{ch02mikolov2013} T. Mikolov, K. Chen, G. Corrado, and J. Dean, ``Efficient estimation of word representations in vector space,'' arXiv:1301.3781, 2013. [Online]. Available: \href{https://arxiv.org/abs/1301.3781}{arxiv.org/abs/1301.3781}
\bibitem{ch02vaswani2017} A. Vaswani et al., ``Attention is all you need,'' arXiv:1706.03762, 2017. [Online]. Available: \href{https://arxiv.org/abs/1706.03762}{arxiv.org/abs/1706.03762}
\bibitem{ch02bert2018} J. Devlin, M.-W. Chang, K. Lee, and K. Toutanova, ``BERT: Pre-training of deep bidirectional transformers for language understanding,'' arXiv:1810.04805, 2018. [Online]. Available: \href{https://arxiv.org/abs/1810.04805}{arxiv.org/abs/1810.04805}
\bibitem{ch02gpt22019} A. Radford et al., ``Language models are unsupervised multitask learners,'' OpenAI, 2019. [Online]. Available: \href{https://cdn.openai.com/better-language-models/language-models.pdf}{cdn.openai.com/better-language-models/language-models.pdf}
\bibitem{ch02gpt32020} T. B. Brown et al., ``Language models are few-shot learners,'' arXiv:2005.14165, 2020. [Online]. Available: \href{https://arxiv.org/abs/2005.14165}{arxiv.org/abs/2005.14165}
\bibitem{ch02scaling2020} J. Kaplan et al., ``Scaling laws for neural language models,'' arXiv:2001.08361, 2020. [Online]. Available: \href{https://arxiv.org/abs/2001.08361}{arxiv.org/abs/2001.08361}
\bibitem{ch02chinchilla2022} J. Hoffmann et al., ``Training compute-optimal large language models,'' arXiv:2203.15556, 2022. [Online]. Available: \href{https://arxiv.org/abs/2203.15556}{arxiv.org/abs/2203.15556}
\bibitem{ch02instructgpt2022} L. Ouyang et al., ``Training language models to follow instructions with human feedback,'' arXiv:2203.02155, 2022. [Online]. Available: \href{https://arxiv.org/abs/2203.02155}{arxiv.org/abs/2203.02155}
\bibitem{ch02lamda2022} R. Thoppilan et al., ``LaMDA: Language models for dialog applications,'' arXiv:2201.08239, 2022. [Online]. Available: \href{https://arxiv.org/abs/2201.08239}{arxiv.org/abs/2201.08239}
\bibitem{ch02chatgpt2022} OpenAI, ``Introducing ChatGPT,'' Nov. 30, 2022. [Online]. Available: \href{https://openai.com/blog/chatgpt/}{openai.com/blog/chatgpt}
\bibitem{ch02rag2020} P. Lewis et al., ``Retrieval-augmented generation for knowledge-intensive NLP tasks,'' arXiv:2005.11401, 2020. [Online]. Available: \href{https://arxiv.org/abs/2005.11401}{arxiv.org/abs/2005.11401}
\bibitem{ch02gpt42023} OpenAI, ``GPT-4 technical report,'' arXiv:2303.08774, 2023. [Online]. Available: \href{https://arxiv.org/abs/2303.08774}{arxiv.org/abs/2303.08774}
\bibitem{ch02constitutional2022} Y. Bai et al., ``Constitutional AI: Harmlessness from AI feedback,'' arXiv:2212.08073, 2022. [Online]. Available: \href{https://arxiv.org/abs/2212.08073}{arxiv.org/abs/2212.08073}
\bibitem{ch02cot2022} J. Wei et al., ``Chain-of-thought prompting elicits reasoning in large language models,'' arXiv:2201.11903, 2022. [Online]. Available: \href{https://arxiv.org/abs/2201.11903}{arxiv.org/abs/2201.11903}
\bibitem{ch02react2022} S. Yao et al., ``ReAct: Synergizing reasoning and acting in language models,'' arXiv:2210.03629, 2022. [Online]. Available: \href{https://arxiv.org/abs/2210.03629}{arxiv.org/abs/2210.03629}
\bibitem{ch02toolformer2023} T. Schick et al., ``Toolformer: Language models can teach themselves to use tools,'' arXiv:2302.04761, 2023. [Online]. Available: \href{https://arxiv.org/abs/2302.04761}{arxiv.org/abs/2302.04761}
\bibitem{ch02llama2023} H. Touvron et al., ``LLaMA: Open and efficient foundation language models,'' arXiv:2302.13971, 2023. [Online]. Available: \href{https://arxiv.org/abs/2302.13971}{arxiv.org/abs/2302.13971}
\bibitem{ch02llama312024} Meta AI, ``Introducing Llama 3.1: Our most capable models to date,'' Jul. 23, 2024. [Online]. Available: \href{https://ai.meta.com/blog/meta-llama-3-1/}{ai.meta.com/blog/meta-llama-3-1}
\bibitem{ch02cua2025} OpenAI, ``Computer-Using Agent,'' Jan. 23, 2025. [Online]. Available: \href{https://openai.com/index/computer-using-agent/}{openai.com/index/computer-using-agent}
\bibitem{ch02openai2025o3} OpenAI, ``Introducing OpenAI o3 and o4-mini,'' Apr. 16, 2025. [Online]. Available: \href{https://openai.com/index/introducing-o3-and-o4-mini/}{openai.com/index/introducing-o3-and-o4-mini}
\bibitem{ch02deepseekr12025} D. Guo et al., ``DeepSeek-R1: Incentivizing reasoning capability in LLMs via reinforcement learning,'' arXiv:2501.12948, 2025. [Online]. Available: \href{https://arxiv.org/abs/2501.12948}{arxiv.org/abs/2501.12948}
\bibitem{ch02gpt552026} OpenAI, ``Introducing GPT-5.5,'' Apr. 23, 2026. [Online]. Available: \href{https://openai.com/index/introducing-gpt-5-5/}{openai.com/index/introducing-gpt-5-5}
\end{thebibliography}
